\documentclass[11pt]{article}

\usepackage[margin=1in]{geometry}
\usepackage[T1]{fontenc}
\usepackage[utf8]{inputenc}
\usepackage{lmodern}
\PassOptionsToPackage{hyphens}{url}
\usepackage{xurl}
\usepackage[hidelinks]{hyperref}
\usepackage{booktabs}
\usepackage{tabularx}
\usepackage{array}
\usepackage{enumitem}
\usepackage{titlesec}
\usepackage{fancyhdr}
\usepackage{setspace}

\setlength{\parindent}{0pt}
\setlength{\parskip}{0.7em}
\setlength{\emergencystretch}{3em}
\raggedbottom
\setlength{\headheight}{14pt}
\pagestyle{fancy}
\fancyhf{}
\lhead{ISO/IEC/IEEE 42010:2022 and TOGAF}
\rhead{Handout}
\cfoot{\thepage}

\titleformat{\section}{\large\bfseries}{}{0pt}{}
\titleformat{\subsection}{\normalsize\bfseries}{}{0pt}{}

\title{\textbf{ISO/IEC/IEEE 42010:2022 and TOGAF}\\[0.35em]
\large A concise handout on their relationship in enterprise architecture}
\author{}
\date{May 7, 2026}

\begin{document}
\maketitle
\thispagestyle{fancy}

\begin{spacing}{1.05}

\textbf{Executive takeaway.} ISO/IEC/IEEE 42010:2022 and TOGAF are complementary, not competing. ISO 42010 defines the requirements for a conformant architecture description; TOGAF provides an enterprise architecture methodology and framework that teams can use to develop and govern architecture work in practice. In short: ISO 42010 is primarily about \emph{what a sound architecture description must contain}; TOGAF is primarily about \emph{how an enterprise architecture practice can organize and execute architecture work}.\cite{iso42010,togaf,togafadm}

\section{What ISO/IEC/IEEE 42010:2022 does}

ISO/IEC/IEEE 42010:2022 specifies requirements for the \emph{structure and expression} of an architecture description (AD) for a broad range of entities, including software, systems, enterprises, systems of systems, product lines, service lines, technologies, and business domains. The standard distinguishes the architecture of an entity of interest from the architecture description that expresses that architecture. It also specifies requirements for related concepts such as architecture description frameworks, architecture description languages, architecture viewpoints, and model kinds.\cite{iso42010,ieee42010}

Just as important, ISO 42010 states what it does \emph{not} specify. It does not prescribe the processes, architecting methods, models, notations, techniques, tools, format, or media by which an architecture description is created or managed. That limitation is central to understanding its relationship to TOGAF.\cite{iso42010}

\section{What TOGAF does}

The Open Group describes TOGAF as an enterprise architecture methodology and framework. The current standard reference point is the \emph{TOGAF Standard, 10th Edition}. The Open Group also maintains a certification portfolio that still includes both Version 9.2 and 10th Edition, which explains why both versions still appear in practice and in training contexts. For current comparison work, however, 10th Edition is the correct baseline.\cite{togaf,togaf10,togafcert}

TOGAF is not limited to describing the contents of architecture documents. It provides an operating structure for enterprise architecture work, including the Architecture Development Method (ADM), which The Open Group describes as the process for moving from foundation architecture material to an organization-specific architecture that addresses business requirements.\cite{togafadm}

\section{The relationship in one sentence}

\textbf{ISO 42010 is the architecture-description standard; TOGAF is the enterprise-architecture framework and method that can be used to produce architecture work products consistent with that standard.}\cite{iso42010,togaf,togafadm}

\section{Comparison at a glance}

\renewcommand{\arraystretch}{1.2}
\begin{tabularx}{\textwidth}{@{}>{\bfseries\raggedright\arraybackslash}p{0.22\textwidth}>{\raggedright\arraybackslash}X>{\raggedright\arraybackslash}X@{}}
\toprule
Dimension & ISO/IEC/IEEE 42010:2022 & TOGAF \\
\midrule
Primary role & Standard for architecture descriptions & Enterprise architecture methodology and framework \\
Core question answered & What must be present in a sound architecture description? & How should an enterprise architecture practice organize, develop, and govern architecture work? \\
Scope emphasis & Structure, expression, viewpoints, model kinds, conformance & Method, practice, guidance, ADM, and supporting architecture content \\
Prescriptiveness & Descriptive about architecture-description requirements & Procedural and practice-oriented for enterprise architecture work \\
What it does not try to be & A full architecting method or toolset & A formal international standard for architecture-description conformance \\
Best use in practice & Quality and conformance lens for architecture descriptions & Delivery and operating model for enterprise architecture \\
\bottomrule
\end{tabularx}

\section{Shared and compatible concepts}

ISO 42010 is built on a conceptual model for architecture description. The standard and its supporting conceptual model emphasize concepts such as stakeholders, concerns, viewpoints, views, and the relationships among them. The architecture-description guidance associated with ISO 42010 also highlights the importance of recording architecture rationale for key decisions. These concepts are compatible with TOGAF-based enterprise architecture practice and provide a disciplined way to test whether architecture artifacts are actually addressing stakeholder concerns in a traceable manner.\cite{conceptualmodel,ads}

\section{How to use them together in an enterprise architecture practice}

A practical way to combine the two is:

\begin{enumerate}[leftmargin=1.5em]
    \item Use TOGAF as the \textbf{practice framework}: organize the architecture function, run the ADM or equivalent enterprise process, assign roles, and govern architecture work products.\cite{togaf,togafadm}
    \item Use ISO 42010 as the \textbf{quality and conformance lens}: check whether each architecture description identifies stakeholders, captures concerns, defines or applies appropriate viewpoints, and produces views and related model kinds that are fit for purpose.\cite{iso42010,conceptualmodel}
    \item Treat architecture documentation as a governed deliverable: the team may use TOGAF to generate architecture content, but the resulting architecture description should still be reviewed against ISO 42010 requirements.\cite{iso42010,togaf}
\end{enumerate}

This combined approach avoids a common failure mode in enterprise architecture: having a method without disciplined architecture-description quality, or having a documentation standard without an effective operating method.

\section{What not to overstate}

Two cautions improve accuracy:

\begin{itemize}[leftmargin=1.5em]
    \item It is safer to say that TOGAF is \textbf{compatible with} ISO 42010 than to claim, without a specific Open Group source in hand, that TOGAF was formally revised \emph{to align with} every detail of ISO 42010:2022.
    \item It is also more precise to compare current ISO 42010 work against \textbf{TOGAF Standard, 10th Edition}, while acknowledging that Version 9.2 remains visible in certification and legacy practice.\cite{togaf10,togafcert}
\end{itemize}

\section{Bottom line}

For enterprise architecture teams, the clean mental model is straightforward: \textbf{use TOGAF to run the practice, and use ISO 42010 to judge the adequacy of the architecture description}. That division of labor preserves the strengths of both: TOGAF gives the team a repeatable way to conduct architecture work, while ISO 42010 provides a rigorous, standard-based basis for reviewing what the architecture description actually says and how well it addresses stakeholder concerns.\cite{iso42010,togaf,togafadm,conceptualmodel}

\section*{References}
\begin{thebibliography}{9}

\bibitem{iso42010}
ISO,
\emph{ISO/IEC/IEEE 42010:2022 - Software, systems and enterprise --- Architecture description},
2022.
\newline\url{https://www.iso.org/standard/74393.html}

\bibitem{ieee42010}
IEEE Standards Association,
\emph{IEEE/ISO/IEC 42010-2022},
2022.
\newline\url{https://standards.ieee.org/ieee/42010/6846/}

\bibitem{conceptualmodel}
ISO 42010 Architecture Description Working Group,
\emph{ISO/IEC/IEEE 42010: Conceptual Model}.
\newline\url{https://www.iso-architecture.org/42010/cm/}

\bibitem{ads}
ISO 42010 Architecture Description Working Group,
\emph{ISO/IEC/IEEE 42010: Architecture Descriptions}.
\newline\url{https://www.iso-architecture.org/ieee-1471/ads/}

\bibitem{togaf}
The Open Group,
\emph{The TOGAF Standard}.
\newline\url{https://www.opengroup.org/togaf}

\bibitem{togaf10}
The Open Group,
\emph{TOGAF Standard, 10th Edition Downloads}.
\newline\url{https://www.opengroup.org/togaf-standard-10th-edition-downloads}

\bibitem{togafcert}
The Open Group,
\emph{TOGAF Certification Portfolio}.
\newline\url{https://www.opengroup.org/certifications/togaf-certification-portfolio}

\bibitem{togafadm}
The Open Group,
\emph{Introduction to the Architecture Development Method (ADM)}.
\newline\url{https://www.opengroup.org/public/arch/p2/p2_intro.htm}

\end{thebibliography}

\end{spacing}
\end{document}



